% Figure 1: architecture. Drawn in TikZ so it needs no external asset.
% Mirrors the reference papers' opening figure: the two context policies side
% by side, plus what each does to prompt growth.
%
% This figure contains NO measured numbers - it is a diagram of mechanism.
% The measured story is Figures 2-4.
%
% Geometry note: the whole picture is scaled to \textwidth at the end. Keep
% new elements inside x in [-3.6, 11.4] or they will push past the margin.

\begin{figure*}[t]
\centering
\resizebox{\textwidth}{!}{%
\begin{tikzpicture}[
  font=\small,
  box/.style={draw, rounded corners=2pt, minimum width=3.4cm,
              minimum height=0.6cm, align=center, inner sep=2pt},
  hist/.style={box, fill=red!8},
  state/.style={box, fill=blue!10, draw=blue!60, thick},
  plain/.style={box, fill=black!3},
  llm/.style={box, fill=black!8, minimum height=0.85cm, thick},
  grp/.style={draw, dashed, rounded corners=3pt, inner sep=7pt},
  arr/.style={-latex, thick},
]

\node[font=\bfseries] at (0,4.55) {Рантайм транскрипта (контроль)};
\node[plain] (p1)  at (0, 3.05) {спецификация задачи $P$};
\node[hist]  (h1)  at (0, 2.35) {история наблюдение--действие};
\node[hist]  (h2)  at (0, 1.65) {предыдущее рассуждение $R_{<t}$};
\node[plain] (o1)  at (0, 0.95) {последнее наблюдение $o_t$};
\node[grp, fit=(p1)(h1)(h2)(o1)] (g1) {};
\node[above=2pt of g1, font=\bfseries\footnotesize, text=red!65!black]
  {контекст $|C_t| = \mathcal{O}(t)$};
\node[llm] (m1) at (0,-0.35) {языковая модель};
\draw[arr] (g1.south) -- (m1.north);
\node[plain, dashed] (a1) at (0,-1.45) {выполнить действие $a_t$};
\draw[arr] (m1) -- (a1);
\draw[arr, red!60!black] (a1.west) -- ++(-1.25,0) |- (h1.west)
  node[pos=0.28, left, font=\scriptsize, text=red!60!black, align=center]
  {история\\растёт};

\begin{scope}[xshift=7.8cm]
\node[font=\bfseries] at (0,4.55) {Рантайм состояния, хвост $k$};
\node[plain] (p2)  at (0, 3.05) {спецификация задачи $P$};
\node[hist]  (t2)  at (0, 2.35) {последние $k$ обменов};
\node[state] (s2)  at (0, 1.65) {состояние выполнения $\Sigma_{t-1}$};
\node[plain] (o2)  at (0, 0.95) {последнее наблюдение $o_t$};
\node[grp, fit=(p2)(t2)(s2)(o2)] (g2) {};
\node[above=2pt of g2, font=\bfseries\footnotesize, text=blue!60!black]
  {контекст $|C_t| = \mathcal{O}(|P|{+}|\Sigma|{+}|o|)$};
\node[right=10pt of t2, font=\scriptsize, text=black!55, align=left]
  {$k{=}0$: пусто};
\node[llm] (m2) at (0,-0.35) {языковая модель};
\draw[arr] (g2.south) -- (m2.north);
\node[plain, dashed] (a2) at (0,-1.45) {выполнить действие $a_t$};
\draw[arr] (m2) -- (a2);
\draw[arr, blue!60!black] (m2.east) -- ++(1.15,0) |- (s2.east)
  node[pos=0.26, right, font=\scriptsize, text=blue!60!black, align=left]
  {проверенный\\патч $\Delta\Sigma_t$};
\node[draw, dashed, text=black!45, minimum width=2.1cm, minimum height=0.55cm,
      rounded corners=2pt, font=\footnotesize] (disc) at (-3.15,-0.35)
      {рассуждение $R_t$};
\draw[arr, dashed, black!45] (m2.west) -- (disc.east);
\node[below=0pt of disc, font=\scriptsize, text=black!45] {отброшено};
\end{scope}

\draw[dashed, black!30] (-3.4,-2.35) -- (11.2,-2.35);

\node[font=\bfseries\small] at (0,-2.75) {рост промпта};
\begin{scope}[yshift=-3.25cm, xshift=-1.9cm]
  \draw[->] (0,0) -- (4.1,0) node[right, font=\scriptsize] {$|C_t|$};
  \draw[->] (0,0) -- (0,-1.15);
  \foreach \y/\w/\lab in {-0.25/0.55/{$t{=}1$}, -0.58/1.7/{$t{=}50$},
                          -0.91/3.35/{$t{=}200$}} {
    \fill[red!35] (0,\y-0.09) rectangle (\w,\y+0.09);
    \node[left, font=\scriptsize] at (0,\y) {\lab};
  }
  \draw[dashed, red] (3.75,0.05) -- (3.75,-1.05);
  \node[font=\scriptsize, text=red, anchor=south] at (3.75,0.05) {окно};
  \node[right, font=\scriptsize, text=red!70!black] at (3.38,-0.91) {растёт};
\end{scope}

\node[font=\bfseries\small] at (7.8,-2.75) {рост промпта};
\begin{scope}[yshift=-3.25cm, xshift=5.9cm]
  \draw[->] (0,0) -- (4.1,0) node[right, font=\scriptsize] {$|C_t|$};
  \draw[->] (0,0) -- (0,-1.15);
  \foreach \y/\w/\lab in {-0.25/0.75/{$t{=}1$}, -0.58/0.8/{$t{=}50$},
                          -0.91/0.85/{$t{=}200$}} {
    \fill[blue!35] (0,\y-0.09) rectangle (\w,\y+0.09);
    \node[left, font=\scriptsize] at (0,\y) {\lab};
  }
  \draw[dashed, red] (3.75,0.05) -- (3.75,-1.05);
  \node[font=\scriptsize, text=red, anchor=south] at (3.75,0.05) {окно};
  \node[right, font=\scriptsize, text=blue!70!black] at (0.9,-0.91) {ограничен};
\end{scope}

\end{tikzpicture}}
\caption{Две политики контекста из Алгоритмов~\ref{alg:transcript}
и~\ref{alg:state}. \textbf{Слева:} рантайм транскрипта дополняет каждый
обмен, так что промпт растёт вместе с эпизодом (Ур.~\ref{eq:ctx-transcript})
и в итоге упирается в окно контекста. \textbf{Справа:} рантайм состояния
несёт проверенный блок состояния и, при $k{=}0$, вообще никакой истории
(Ур.~\ref{eq:ctx-state}); его промпт ограничен независимо от $t$. Оба
отбрасывают промежуточное рассуждение. Ограниченный промпт — это
опубликованная выгода. Чего диаграмма не может показать — это величина,
которую измеряет данная статья: ограничение $|C_t|$ не ограничивает
информацию, требуемую задачей, так что агент заново её получает через
вызовы инструментов, раздувая реализованный горизонт
(Ур.~\ref{eq:horizon-inflation}) без затрат, видимых какой-либо метрике
размера промпта.}
\label{fig:architecture}
\end{figure*}
